\subsection{经期性行为的实操与舒适度：体位、浸染与用品}

前面一节回答了"能不能"以及风险与避孕。这一节处理剩下的问题：\textbf{怎样让它舒服、少狼狈、双方都不别扭}。

\textbf{一、先理解身体在经期有什么不同}

\begin{itemize}
\item \textbf{盆腔充血}：部分女性因此唤起更充分、感受更敏感；也有一部分人感到胀、钝痛，性欲反而降低。两种反应都常见，且同一个人在不同经期也可能不同；
\item \textbf{宫颈更敏感}：经期宫颈位置与敏感度有变化，深插入容易引发不适（尤其后入式与深插入姿势）。若出现这种不适，缩短插入深度、放慢节奏或更换体位通常即可改善，不必硬撑；
\item \textbf{润滑情况改变}：经血本身有润滑作用，但也可能对部分人的外阴皮肤形成刺激。使用安全套时，经血与乳胶之间的摩擦感与平时不同，可适当补充\textbf{水性润滑剂}（不要用油性润滑剂，会损坏乳胶安全套）；
\item \textbf{腹胀与腹压}：腹部受压或需要核心用力维持的姿势，在腹胀、痛经明显的日子会加重不适。相对友好的选择是\textbf{侧卧位}（腹部不受压、插入较浅）、\textbf{女上位}（自己控制深度与节奏）、\textbf{双腿抬高的传教士位}（承重由对方支撑）。
\end{itemize}

\textbf{二、把"麻烦"降到可接受（可执行的准备）}

\begin{itemize}
\item \textbf{事前}：床上铺深色毛巾或一条旧浴巾；或直接选择浴室、淋浴环境，把"清理"这件事当场解决。若经量较大，可优先选择经量较少的时段（多在经期第二、三天之后）；
\item \textbf{事后}：及时排尿（有助于预防尿路感染）、用温水清洗外阴，\textbf{不要冲洗阴道内部}；
\item \textbf{用品更换}：结束后按常规更换卫生用品，保持干燥清洁。
\end{itemize}

\textbf{三、用品的"性爱友好度"}

不同用品在这件事上的适用性差别很大，值得单独说明：

\begin{itemize}
\item \textbf{月经杯}：\textbf{软质杯}（一次性软杯或杯壁较软的硅胶杯）通常可以在性生活时戴用，同时起到收集经血与一定屏障的作用——这是经期性行为里最实用、也最常被忽略的一条；\textbf{硬质杯}（杯壁较硬、带明显排气孔结构）建议取下，因为可能被顶移位、造成不适或取出困难；杯的尾柄若被触碰不适，可在说明允许范围内剪短；
\item \textbf{避孕膜（子宫帽）}：本身就是屏障避孕工具，同时也能部分承接经血；
\item \textbf{卫生棉条}：\textbf{必须取出后}再进行——被推入深处会导致取出困难、增加感染与中毒性休克综合征风险；
\item \textbf{卫生巾}：事前取下。也可使用一次性垫巾，事后直接丢弃，减少清洗负担；
\item \textbf{不论使用哪种用品}：安全套的意义不因用品而改变——它同时降低性传播感染与盆腔感染风险。
\end{itemize}

\textbf{四、羞耻感与沟通}

\begin{itemize}
\item 关于月经的禁忌在几乎所有文化中都能找到，但\textbf{它的来源是文化禁忌，不是卫生问题}。从小被教育"经期不能做某事"，并不意味着这件事本身有什么错——也不意味着你必须照旧执行这套规则；
\item \textbf{开口可以很简单}：直接、简短地说明意愿与边界即可。例如"今天可以，但我想铺条毛巾""今天肚子胀，我们换个别的方式""今天我不想做"。三句话都不需要额外解释；
\item \textbf{也要允许"想在经期完全不做"这个选项}：不做的理由同样不需要提供；
\item \textbf{不必表演接受}：为了证明自己"不封建"而勉强进行，和为了守规矩而勉强拒绝，是同一类伤害。判断标准只有一个——你自己是否愿意。
\end{itemize}

\begin{tcolorbox}[colback=pink!5!white,colframe=red!50!black,title=贴心小叮咛]
经期性行为的目标不是"克服羞耻完成一次"，而是把它变成一件和平时一样\textbf{可以被商量}的事。商量得顺，做与不做都不会留下负担。
\end{tcolorbox}

